%======================================================================
% صفحه‌ی عنوان
%======================================================================
\begin{titlepage}
\thispagestyle{empty}
\begin{center}

% ---- بسم الله ----
{\Large بسم الله الرحمن الرحیم}

\vspace{1.2cm}

% ---- لوگو دانشگاه ----
% \includegraphics[width=3.5cm]{figures/aut_logo.png}\\[0.6cm]

% ---- نام دانشگاه ----
{\LARGE\bfseries دانشگاه صنعتی امیرکبیر}\\[0.3cm]
{\large (پلی‌تکنیک تهران)}\\[0.2cm]
{\large دانشکده‌ی مهندسی کامپیوتر}

\vspace{0.8cm}
\hrule height 1pt
\vspace{0.5cm}

% ---- نوع مدرک ----
{\large پایان‌نامه‌ی کارشناسی}\\[0.3cm]
{\normalsize در رشته‌ی مهندسی کامپیوتر}

\vspace{0.8cm}

% ---- عنوان ----
{\LARGE\bfseries
طراحی و پیاده‌سازی سامانه‌ی واسط برای\\[0.3cm]
یکپارچه‌سازی سرویس‌های هوش مصنوعی
}

\vspace{0.4cm}
{\large\lr{Design and Implementation of a Unified Backend Gateway\\[0.2cm]
for AI Service Integration}}

\vspace{0.8cm}
\hrule height 0.5pt
\vspace{0.8cm}

% ---- مشخصات ----
\begin{tabular}{r@{\hspace{0.5cm}}l}
\textbf{ارائه‌دهنده:}    & مصطفی رحمتی \\[0.4cm]
\textbf{شماره دانشجویی:} & \lr{9931069} \\[0.4cm]
\textbf{استاد راهنما:}   & \textbf{[نام استاد راهنما]} \\[0.4cm]
\end{tabular}

\vspace{1.2cm}

% ---- تاریخ ----
{\large نیمسال دوم سال تحصیلی ۱۴۰۴–۱۴۰۵}

\vspace{0.5cm}
{\normalsize بهار ۱۴۰۵}

\end{center}
\end{titlepage}
\cleardoublepage
